\documentclass{article}
\usepackage[utf8]{inputenc}
\usepackage{amsmath}
\usepackage{amssymb}
\usepackage{graphicx}
\usepackage{tabularx}
\usepackage{booktabs}
\usepackage[english]{babel}
\usepackage{setspace}
\usepackage{geometry}
\usepackage{float}
\usepackage{url}
\usepackage{listings}
\usepackage{xcolor}
\usepackage{longtable}
\usepackage[justification=centering]{caption}

\lstset{
    basicstyle=\ttfamily\small,
    backgroundcolor=\color{gray!10},
    frame=single,
    breaklines=true,
    columns=fullflexible
}

% APA style formatting
\geometry{a4paper, margin=1in}
\setstretch{1.5}
\captionsetup[table]{labelfont=bf, justification=centering, singlelinecheck=false}
\captionsetup[longtable]{labelfont=bf, justification=centering, singlelinecheck=false}
\captionsetup[figure]{font=small, labelfont=bf, justification=centering, singlelinecheck=false}

\title{ Machine Learning for Customer Churn Prediction}
\author{Nelson David Posso Suarez \\ Javier Alejandro Penagos Hernandez \\ Nicolas Castro Rivera }
\date{}

\begin{document}

\maketitle


\section{Executive Summary}
Unsupervised learning provides techiniques useful for different problems that complement the projects. The re-creation of
existing data, the anomaly detection or generate new data similar to existing data. Can be help in the domain of Churn Prediction,
in special for working with a dataset like the KKBox selected for this project. 

\section{Motivation}
The contribution of this workshop to the project, is implement unsupervised learning for strategy parts in the Customer Churn
Prediction system. There is a problem with the unbalanced dataset, then the experiment will consist in explore the AE and VAE
for show how good reply existing customers and then implement for generate customer churns. 
\\Purpose of this new customers generated with real customer churns, is generate more data and get a better performace, that 
the models of w1 got. The experiment will conclude if this whether or not this is feisible depends on the results. 



\section{AE and VAE}
Autoencoder and variational autoencoder are unsupervised neural network arquitectures

\subsection{Data and preprocessing}

The autoencoders reuse the tabular export built in Workshop~1 (version~5), where the
\textit{members}, \textit{transactions} and \textit{user\_logs} tables are condensed into a
single record of 59 features per customer (56 numerical and 3 categorical). The February
cohort (992{,}931 customers, 6.39\% churn) is used for development and the March cohort
(970{,}960 customers, 8.99\% churn) as the held-out test set; 881{,}701 customers appear in
both months.

The Workshop~1 split (stratified and hash-based on \texttt{msno}) is reused so that the
same customers are used across workshops. The churn label is never fed to the networks; it
is only used to select the training population. The AE and the VAE are trained on the
697{,}095 non-churn customers of the \textit{fit} partition and validated on the 116{,}183
non-churn customers of the \textit{stop} partition, while the full March cohort, including
churners, is used for testing.

A preprocessor is fitted on the training population only. Numerical features are imputed
with the median, skewed features are transformed with a signed $\log(1+|x|)$, clipped to the
0.1\% and 99.9\% quantiles and standardized; one constant column was removed, leaving 55
numerical inputs. The categorical features (\texttt{city}, \texttt{registered\_via},
\texttt{last\_payment\_method\_id}) are mapped to vocabularies of 22, 6 and 34 levels and
one-hot encoded inside the network, giving an input dimension of 117. An inverse transform
maps reconstructions back to the original 59 columns. The AE and the VAE share the same
fitted preprocessor.

\subsection{AE Architecture}

\subsection{AE Train}

\subsection{VAE Architecture}

\subsection{VAE Train}

\section{Reconstruction quality}

\section{Anomaly detection}

\section{Latent space visualization}

\section{Generate synthetic data}

\section{Analysis}

\begin{thebibliography}{9}
\bibitem{rizvi2018kkbox}
Rizvi, A. (2018). \textit{WSDM -- KKBox's Churn Prediction Challenge}. Medium. Retrieved from \url{https://asad-99rizvi.medium.com/wsdm-kkboxs-churn-prediction-challenge-d7c9cfa21cbd}
\end{thebibliography}

\end{document}
